% Vietnamese translation for OI Public Library
% Source-File: IOI中国国家候选队论文/2003/侯启明/侯启明.ppt
\documentclass[11pt]{article}
\usepackage{oipl-vn}

\title{Bản trình chiếu: Ứng dụng lý thuyết thông tin trong thi tin học}
\author{Hou Qiming}
\date{2003}

\begin{document}
\maketitle

\origtitle{Slide companion for information-theory applications}
\sourcepath{IOI China candidate team papers / 2003 / Hou Qiming / Hou Qiming.ppt}

\section{Phạm vi bản dịch}

Bản trình chiếu này đi cùng bài viết về entropy và cận dưới thông tin. Phần
chữ khôi phục được chứa các công thức \(H(x)\), các phân phối xác suất, và
các ví dụ \textit{Rods} của IOI 2002 cùng bài \textit{Coins}.

\section{Entropy}

Với biến ngẫu nhiên có phân phối \(p_0,p_1,\ldots,p_n\), lượng thông tin được
đo bằng
\[
  H(x)=-\sum_i p_i\log p_i .
\]
Nếu các khả năng đều nhau, entropy đạt cực đại; chẳng hạn tám khả năng đều
nhau cho \(H(x)=\log 8\).

\section{Ứng dụng}

Trong \textit{Rods}, slide dùng số lượng hình chữ nhật khả dĩ
\(\operatorname{rect}(a,b,c,d)\) để ước lượng không gian trạng thái, rồi suy ra
cận dưới kiểu \(6\log_2 n+C\). Với các bài cân đồng xu, đầu ra mỗi phép cân có
tối đa ba kết quả, nên cận dưới tự nhiên là \(\log_3 n\).

\section{Ý nghĩa}

Lý thuyết thông tin cung cấp ngôn ngữ thống nhất để chứng minh vì sao một số
lời giải không thể dùng ít truy vấn hơn. Bản trình chiếu nhấn mạnh vai trò
của việc đếm trạng thái và chọn phép đo có entropy lớn.

\end{document}
